% SHARD-ID: AQM-80-QSA-FINITE-RING-EXAMPLE-MATRIX
% SHARD-TITLE: Finite Ring Example Matrix
% SHARD-SUMMARY: Builds a gap-aware finite-ring matrix for QSA Step 16 without claiming a populated finite-ring database.
% SHARD-SUMMARY: Separates point/residue, cotangent-first, field-label, and nilpotent-sensitive workflows row by row.
% SHARD-SUMMARY: Carries explicit available, blocked, unknown, not_applicable, and degenerate markers to later QSA steps.
% SHARD-KEYWORDS: quantum-system association, finite rings, example matrix, residue fields, cotangent, nilpotents, database gaps

\section{Finite Ring Example Matrix}
\label{sec:qsa-finite-ring-example-matrix}

\subsection{Status and Local Anchors}

\textbf{Status: Synthesis matrix.}  This is QSA Step 16 from
\path{docs/quantum_system_associations/PLAN.md}.  It adds no source,
convention entry, script, run bundle, generated data, populated database row,
or finite-ring classification claim.  Its purpose is to put the standard
finite-ring examples into one gap-aware comparison table before the dedicated
residue-site and nilpotent-sensitive shards.

The local anchors are AQM-23, AQM-34--AQM-37, AQM-40, AQM-58, and
AQM-76--AQM-79.  The finite-ring database boundary is the PRD
\path{docs/finite_commutative_ring_database_prd.md}, the implementation and
population plans, convention \texttt{(an)}, the schema
\path{data/SCHEMA.md}, and
\path{runs/2026-05-26-finite-ring-database/README.md}.  That run bundle
records a schema-only SQLite smoke build with one \texttt{build\_run} row and
zero \texttt{ring} or \texttt{presentation} rows.  Its CSV files are
in-memory MVP review exports; they do not populate or audit
\path{finite_rings.sqlite}.

The open blockers are part of the table.  General finite-ring residue
decomposition and deterministic site ordering remain blocked by
\texttt{aqm-qsa-frres01}.  Database-backed intrinsic or relative cotangent
columns remain blocked by \texttt{aqm-qsa-frcot01}.  Two-direction local Artin
standardization remains blocked by \texttt{aqm-qsa-artin01} except for the
local cotangent derivation in AQM-77.  Quotient-collapse examples beyond the
existing \(\mathbb Z/6\mathbb Z\cong\mathbb F_2\times\mathbb F_3\) anchor
remain blocked by \texttt{aqm-qsa-collapse01}.

\subsection{Status Tokens}

The table uses the report-planning labels of convention \texttt{(au)}.  They
are not CSV sentinel rows, do not change \path{data/SCHEMA.md}, and do not
alter the \texttt{\#}-prefixed sentinel comments in the finite-ring review
exports.  Here \texttt{available} means that a named local shard or review
export supplies the row; \texttt{blocked} means that a named source,
convention, derivation, helper, or bead is missing; \texttt{unknown} means the
matrix has no local determination yet; \texttt{not\_applicable} means the
workflow is outside the row's stated semantics; and \texttt{degenerate} marks
a locally recorded collapse or forgotten-structure case.

\subsection{Points, Residues, and Cotangent Data}

\begin{center}
\scriptsize\setlength{\tabcolsep}{3pt}
\begin{tabular}{p{0.20\linewidth}p{0.20\linewidth}p{0.22\linewidth}p{0.30\linewidth}}
\toprule
Ring row & Points & Residue data & Cotangent-first data \\
\midrule
\(\mathbb F_3\) &
\texttt{available}: one \(\Spec\)-point by AQM-22 and convention
\texttt{(u)}. &
\texttt{available}: residue field \(\mathbb F_3\), one qutrit site. &
\texttt{available}: AQM-77 gives
\(C_x^{\rm loc}=C_x^{\rm rel}=0\); the cotangent-first carrier is
one-dimensional. \\
\addlinespace
\(\mathbb F_3\times\mathbb F_3\) &
\texttt{available}: two discrete product-spectrum points by AQM-40 and
convention \texttt{(ad)}. &
\texttt{available}: two \(\mathbb F_3\) residue sites, hence two qutrit
factors in the residue workflow. &
\texttt{degenerate}: AQM-58 and AQM-77 give
\(\Omega^1_{\mathbb F_3^2/\mathbb F_3}=0\); the cotangent-first labels vanish
at both points. \\
\addlinespace
\(\mathbb Z/6\mathbb Z\) &
\texttt{available}: AQM-23 derives the two points \((2)\) and \((3)\). &
\texttt{available}: residues \(\mathbb F_2\) and \(\mathbb F_3\), with label
group \(\mathbb F_2^2\oplus\mathbb F_3^2\). &
\texttt{unknown}: no local cotangent comparison row is recorded here;
database-backed cotangent data are blocked by \texttt{aqm-qsa-frcot01}. \\
\addlinespace
\(\mathbb F_3[\epsilon]/(\epsilon^2)\) &
\texttt{available}: the polynomial-quotient/local-Artin point is the unique
closed point in AQM-34 and AQM-77. &
\texttt{available} locally as one \(\mathbb F_3\) residue site; the database
residue helper row is still \texttt{blocked} by \texttt{aqm-qsa-frres01}. &
\texttt{available}: AQM-77 gives one intrinsic and relative
\(\mathbb F_3\)-cotangent direction over \(\mathbb F_3\), carrier dimension
\(3\). \\
\addlinespace
\(\mathbb F_3[t]/(t^3)\) &
\texttt{available}: one closed point, a triple thickening of \(t=0\), by
AQM-34 and AQM-77. &
\texttt{available}: one \(\mathbb F_3\) residue site; the reduced residue
layer is one qutrit. &
\texttt{available}: AQM-77 gives one first-order cotangent direction in
characteristic \(3\); this does not record the full length-three thickening. \\
\addlinespace
\(\mathbb F_3[t]/(t^n-1)\) &
\texttt{available}: AQM-35 uses the irreducible factors of the squarefree
part \(t^m-1\), \(m=n/3^{v_3(n)}\). &
\texttt{available}: each factor \(\pi\) contributes
\(\mathbb F_3[t]/(\pi)\); repeated factors do not add residue sites. &
\texttt{unknown}: no uniform cotangent table is claimed for the family in
this matrix. \\
\addlinespace
\(\mathbb F_3[u,v]/(u^2,v^2)\) &
\texttt{available} only as the local AQM-77 cotangent example at
\((u,v)\); standard Artin-row status is \texttt{blocked}. &
\texttt{available} locally as residue \(\mathbb F_3\) in AQM-77; no
database residue row or standardized finite-ring catalogue row is claimed. &
\texttt{available}: over \(\mathbb F_3\), AQM-77 gives two cotangent
directions and carrier dimension \(9\); relative to
\(\mathbb F_3[u]/(u^2)\), it gives one direction. \\
\addlinespace
\(\mathbb Z/4\mathbb Z\) &
\texttt{available} only for the AQM-77 local cotangent comparison at
\((2)\). &
\texttt{available} locally as residue \(\mathbb F_2\) in AQM-77; the database
residue helper row is \texttt{blocked}. &
\texttt{available}: AQM-77 gives
\(C_x^{\rm loc}\cong\mathbb F_2\) but
\(C_x^{\rm rel}(\mathbb Z)=0\). \\
\addlinespace
\(\mathbb F_2[x]/(x^2+x)\) candidate &
\texttt{blocked}: possible quotient-collapse row only. &
\texttt{blocked}: do not import product points before
\texttt{aqm-qsa-collapse01} supplies a local derivation or certificate. &
\texttt{not\_applicable}: no cotangent comparison is made before the row is
accepted as checked. \\
\bottomrule
\end{tabular}
\end{center}

\subsection{Field Labels, Nilpotents, and Run Artifacts}

\begin{center}
\scriptsize\setlength{\tabcolsep}{3pt}
\begin{tabular}{p{0.20\linewidth}p{0.25\linewidth}p{0.25\linewidth}p{0.22\linewidth}}
\toprule
Ring row & Field-label data & Nilpotent-sensitive data & Run artifact or gap marker \\
\midrule
\(\mathbb F_3\) &
\texttt{available}: one-site finite phase labels; the underlying-set
\(\mathbb F_3\) all-map system is a separate AQM-22/AQM-71 object. &
\texttt{not\_applicable} as a nilpotent-thickening example; the database
thickened helper still records a blocked policy row. &
Finite-ring review CSV: residue row \texttt{available}; SQLite is
schema-only, not populated. \\
\addlinespace
\(\mathbb F_3\times\mathbb F_3\) &
\texttt{available}: AQM-40 gives the product-field residue label net; bare
two-point all-map rows are separate finite-set examples. &
\texttt{not\_applicable} as a nilpotent-thickening example; no thickened
product policy is imported. &
Finite-ring review CSV: residue row \texttt{available}; thickened row
\texttt{blocked}; no database audit claim. \\
\addlinespace
\(\mathbb Z/6\mathbb Z\) &
\texttt{available}: mixed-residue label direct sum in AQM-23; a single
\(\mathbb F_3\)-linear field-label reading is \texttt{not\_applicable}. &
\texttt{blocked}: no certified generating character or thickened policy is
recorded for this row. &
\texttt{degenerate}: the review export records the
\(\mathbb F_2\times\mathbb F_3\) presentation merged to the
\(\mathbb Z/6\mathbb Z\) representative; no extra collapse example is used. \\
\addlinespace
\(\mathbb F_3[\epsilon]/(\epsilon^2)\) &
\texttt{available}: reduced coordinate-ring field profiles forget
\(\epsilon\), as in AQM-34. &
\texttt{available}: AQM-36 gives the top-coefficient Artin-Weyl layer on
\(\ell^2(R)\), dimension \(9\). &
Finite-ring review CSV: residue row \texttt{blocked}, thickened row
\texttt{available}; SQLite has no ring rows. \\
\addlinespace
\(\mathbb F_3[t]/(t^3)\) &
\texttt{available}: reduced field profiles factor through
\(\mathbb F_3\); nilpotent \(t\)-classes vanish at the residue site. &
\texttt{available}: AQM-36/AQM-37 give \(\ell^2(\mathbb F_3[t]/(t^3))\) of
dimension \(27\). &
Local report derivation only; no MVP database presentation or populated
SQLite row is claimed. \\
\addlinespace
\(\mathbb F_3[t]/(t^n-1)\) &
\texttt{available}: AQM-35 field labels depend on the squarefree residue
part in the reduced layer. &
\texttt{available}: AQM-37 gives thickened dimension \(3^n\) and reduced
dimension \(3^m\). &
Local report family only; database population and audit remain pending. \\
\addlinespace
\(\mathbb F_3[u,v]/(u^2,v^2)\) &
\texttt{unknown}: no field-label or residue-site matrix row beyond the
AQM-77 cotangent derivation is standardized. &
\texttt{blocked}: two-direction Artin standardization and any thickened
Weyl row remain under \texttt{aqm-qsa-artin01}. &
No run artifact; use AQM-77 only for cotangent data. \\
\addlinespace
\(\mathbb Z/4\mathbb Z\) &
\texttt{unknown}: no field-label row is standardized in this matrix. &
\texttt{blocked}: finite-ring review CSV records missing certified
generating character; no thickened claim is imported. &
Finite-ring review CSV: residue and thickened rows \texttt{blocked}; AQM-77
only supports the cotangent comparison. \\
\addlinespace
\(\mathbb F_2[x]/(x^2+x)\) candidate &
\texttt{blocked}: possible \texttt{degenerate} quotient-collapse case is not
asserted as checked. &
\texttt{unknown}: no nilpotent-sensitive or reduced-row conclusion is used. &
Blocked by \texttt{aqm-qsa-collapse01}; quotient-helper local-core status is
not a report-level collapse proof. \\
\bottomrule
\end{tabular}
\end{center}

\subsection{Interpretation}

The matrix separates four workflows that are easy to conflate.  The
residue-site workflow uses closed or maximal points and residue fields, as in
AQM-23, AQM-34, AQM-35, and AQM-40.  The cotangent-first workflow uses
intrinsic or relative first-order cotangent data, as in AQM-76 and AQM-77.
The field-label workflow uses finite maps, finite closed supports, or
coordinate-ring profiles, as reviewed in AQM-79.  The nilpotent-sensitive
workflow uses local Artin or Frobenius-ring data, as in AQM-36 and AQM-37.

The table also records the current database boundary.  The 2026-05-26
finite-ring run is useful evidence about the schema-only smoke build and the
in-memory MVP review-export status of selected rows.  It is not a populated
or audited finite-ring database, and no row above treats it as one.

Step 17 must therefore supply certified residue decompositions and site
orderings before extending the residue workflow across finite rings.  Step 18
must keep reduced residue sites separate from nilpotent-sensitive Artin or
Frobenius data, and it must not use the two-direction or quotient-collapse
candidate rows until their blocker beads have local evidence.
